\section{Architecture}
\label{sec:architecture}

The compiler is a short pipeline with three contracts at its joints:
\emph{detection}, the \texttt{DiagramModule}, and the \texttt{Scene}.

\subsection{The pipeline}

\begin{enumerate}
  \item \textbf{Detect} (\texttt{src/frontend/detect.ts}). The source is
        classified by its leading header into an input format
        (\texttt{mermaid} text or \texttt{yaml}/JSON) and a \texttt{DiagramKind}
        (e.g.\ a line starting \texttt{flowchart} $\to$ \texttt{flowchart}; a
        line starting \texttt{avl} $\to$ \texttt{avl}).
  \item \textbf{Dispatch} (\texttt{src/frontend/registry.ts}). The kind selects
        a registered \texttt{DiagramModule}.
  \item \textbf{Parse}. The module turns source into its own IR
        (\texttt{parseMermaid} for text, \texttt{parseYaml} for the structured
        alternate).
  \item \textbf{Layout}. The module places its IR and returns a
        \texttt{LayoutResult} --- a \texttt{Scene} plus a node-anchor registry.
  \item \textbf{Render} (\texttt{src/render/svg.ts}). One renderer turns the
        \texttt{Scene} into SVG; \texttt{resvg} optionally rasterises it to PNG.
\end{enumerate}

Text is the primary surface. The \texttt{yaml}/JSON path is a thin alternate
input (each module's \texttt{parseYaml} accepts a pre-built IR); every example in
this repository is authored as Mermaid-superset text (\texttt{.mmd}).

\subsection{The \texttt{DiagramModule} contract}
\label{sec:diagram-contract}

Every diagram kind implements one interface
(\texttt{src/contracts/diagram.ts}):

\begin{lstlisting}[language=ts]
interface DiagramModule<IR extends BaseIR> {
  parseMermaid(input: string): IR;
  parseYaml(input: string): IR;
  layout(ir: IR, theme: ResolvedTheme): Promise<LayoutResult>;
}
\end{lstlisting}

\texttt{BaseIR} carries only what the pipeline needs uniformly (version,
metadata, optional overlays and per-diagram theme override). Each kind extends it
with its own fields. \texttt{layout} returns:

\begin{lstlisting}[language=ts]
interface LayoutResult {
  scene: Scene;                 // what the renderer draws
  anchors: NodeAnchorRegistry;  // node bounds + cardinal ports
}
\end{lstlisting}

The anchor registry is what makes a diagram \emph{linkable}: it exposes each
node's bounding box (and optional N/S/E/W ports) in the diagram's local
coordinate space, so a composition layer can draw arrows between nodes that live
in different diagrams (\S\ref{sec:composition}).

\subsection{Registries}

Three small in-process registries hold the moving parts: diagram modules
(\texttt{registerDiagram}), renderers (\texttt{registerRenderer}), and edge
routers. Registration is the only wiring needed to add a kind; the rest of the
system discovers it by contract.

\subsection{Build}

The project is a single package at the repository root. Grammars are written in
Peggy (one \texttt{grammar.peggy} per diagram under
\texttt{src/diagrams/<kind>/}); \texttt{scripts/build-grammars.mjs} compiles them
to parsers, and \texttt{tsc} compiles the TypeScript. Tests run under Vitest.
